Ce rapport présente le travail réalisé au sein de l’entreprise \textbf{RUHENNA}, une activité artisanale portée par \textbf{Sumeyra Solak}, dans le cadre d’un stage de 2\textsuperscript{ème} année. L’objectif du projet était de concevoir et développer une plateforme web \textit{full-stack} permettant de renforcer l’image de marque, de valoriser les créations autour de l’art du henné et de simplifier la gestion opérationnelle (catalogue de produits, prise de rendez-vous, blog et formulaire de contact), tout en offrant un \textbf{panneau d’administration} pour centraliser l’édition du contenu et le suivi des demandes.

La solution s’appuie sur une architecture moderne basée sur \textbf{React} et \textbf{Next.js} (App Router), \textbf{Prisma} comme ORM, et une base de données \textbf{PostgreSQL} hébergée sur \textbf{Neon} pour faciliter le déploiement sur \textbf{Vercel}. La gestion des médias s’appuie sur \textbf{Cloudinary} afin d’optimiser le stockage et la diffusion d’images. Une attention particulière a été portée à l’expérience utilisateur (interfaces publiques et administration), à la structuration des données (modèle relationnel, contraintes d’intégrité) et à la maintenabilité (typage TypeScript, validation des entrées, séparation des responsabilités).

Au-delà des aspects techniques, ce stage a constitué une montée en compétences progressive sur React/Next.js et Prisma, depuis un niveau initial limité, à travers un apprentissage guidé par la documentation et des ressources en ligne, puis par la confrontation à des problèmes réels (routage, interaction front/back, disponibilité de créneaux, gestion d’images, déploiement).
